El CVRP ha sido objeto de estudio sistemático en la literatura de investigación
operacional desde su introducción a fines de los años cincuenta \parencite{dantzig1959truck}. Su carácter
$\mathcal{NP}$-hard \parencite{lenstra1981complexity} explica el interés sostenido por
desarrollar métodos exactos capaces de certificar optimalidad sobre instancias de tamaño
no trivial, así como métodos aproximados que escalen a instancias industriales. En este
informe el alcance se restringe deliberadamente a los métodos exactos de programación
matemática, alineados con los contenidos de optimización combinatorial y programación
lineal entera del curso \parencite{tan2021vehicle}.

\subsection{Métodos exactos para el CVRP}

La literatura clasifica los métodos exactos para el CVRP bajo el paraguas
\textit{Branch and X}, donde la base común es el procedimiento de \textit{Branch and
Bound} aplicado sobre formulaciones de programación lineal entera mixta
\parencite{tan2021vehicle}. Las variantes principales se diferencian por la forma en que
fortalecen la relajación lineal, generan variables o exploran el espacio de soluciones.

\subsubsection{Branch and Bound clásico}

El algoritmo de \textit{Branch and Bound} resuelve la relajación lineal del problema
entero, ramifica sobre variables fraccionarias creando subproblemas restringidos, y poda
ramas cuya cota inferior excede la mejor solución entera conocida
\parencite{land1960automatic}. Este esquema constituye el núcleo de los métodos exactos
clásicos cubiertos en el curso de Optimización Combinatorial. En el contexto del CVRP,
la primera implementación efectiva formuló el problema simétrico sobre un grafo no
dirigido y resolvió una relajación inicial con restricciones de grado, agregando
dinámicamente las restricciones de eliminación de subtours y capacidad sólo cuando una
solución entera candidata las incumplía \parencite{laporte1983branch}. Esta estrategia
anticipa, en su lógica esencial, el principio de separación de desigualdades válidas que
caracteriza a los métodos posteriores.

\subsubsection{Branch and Cut}

\textit{Branch and Cut} extiende el esquema anterior incorporando la generación
sistemática de planos de corte que fortalecen la relajación lineal sin alterar el
conjunto de soluciones enteras factibles \parencite{naddef2002branch}. Para el CVRP, las
familias de cortes más relevantes derivan de la estructura poliédrica del problema:
desigualdades de capacidad redondeada, desigualdades de capacidad enmarcada,
desigualdades \textit{comb}, \textit{multistar} y \textit{hypotour}, entre otras
\parencite{naddef1991branch}. El algoritmo de \textcite{lysgaard2004branch} integró estas
familias en un único esquema computacional capaz de resolver por primera vez instancias
del Set B de \textcite{augerat1995} previamente abiertas, como B-n50-k8, B-n66-k9 y
B-n78-k10. \textcite{augerat1995} habían propuesto previamente un Branch and Cut basado
en cortes de capacidad y aportaron el conjunto de instancias A y B que hoy forma parte
estándar de CVRPLIB.

\subsubsection{Branch and Price}

\textit{Branch and Price} aborda el CVRP desde una formulación de partición de conjuntos
en la que cada variable representa una ruta factible completa
\parencite{tan2021vehicle}. Como el número de rutas factibles es exponencial, se recurre
a generación de columnas: el problema maestro se resuelve sobre un subconjunto reducido
de rutas y un subproblema de \textit{pricing} —típicamente un problema de camino más
corto con restricciones de recursos— identifica nuevas rutas con costo reducido
negativo. Este enfoque proporciona cotas inferiores generalmente más fuertes que las
obtenidas por \textit{Branch and Cut} sobre formulaciones de aristas, a cambio de un
subproblema de \textit{pricing} computacionalmente más exigente.

\subsubsection{Branch, Cut and Price}

La combinación de generación de columnas con planos de corte da lugar al esquema
\textit{Branch, Cut and Price} (BCP), considerado el estado del arte actual para el
CVRP. \textcite{fukasawa2006robust} propusieron un BCP robusto que extendió el tamaño
de las instancias resueltas exactamente hasta 135 vértices, integrando cortes sobre la
formulación original con generación de columnas sobre $q$-rutas.
\textcite{pecin2017improved} introdujeron posteriormente cortes \textit{limited-memory}
\textit{rank-1}, que permiten incorporar un mayor número de cortes sin incrementar
proporcionalmente la complejidad del subproblema de \textit{pricing}; esta contribución
permitió resolver instancias clásicas con hasta 420 clientes.

\subsubsection{Programación dinámica}

La programación dinámica fue históricamente uno de los primeros enfoques exactos
explorados para el CVRP. \textcite{christofides1981exact} propusieron algoritmos
basados en relajaciones de árbol de cubrimiento y de camino más corto, que permitieron
resolver instancias de pequeño tamaño y dieron origen al conjunto E utilizado hasta hoy
como banco de pruebas. En la actualidad, la programación dinámica rara vez se utiliza
como método independiente para el CVRP, pero subsiste como componente esencial dentro
del subproblema de \textit{pricing} de los algoritmos \textit{Branch and Price} y
\textit{Branch, Cut and Price}.

\subsection{Posicionamiento del enfoque adoptado}

La progresión histórica desde \textit{Branch and Bound} clásico hacia
\textit{Branch, Cut and Price} refleja un compromiso continuo entre la complejidad de la
formulación, la fortaleza de la relajación lineal y la dificultad del subproblema
auxiliar. Los métodos más sofisticados resuelven instancias considerablemente más
grandes, pero su implementación requiere infraestructura algorítmica especializada
—generación de columnas, separación de cortes poliédricos, etiquetado para
\textit{pricing}— que excede el alcance de un primer informe centrado en programación
matemática clásica.

Por esta razón, el presente trabajo adopta como base metodológica el esquema de
\textcite{laporte1983branch}, complementado con la generación iterativa de
desigualdades de capacidad redondeada (RCI) sobre la formulación de aristas. Este
enfoque permite ilustrar de forma transparente los conceptos centrales que articula el
curso —relajación lineal, integralidad, cotas inferiores, ramificación, poda y
separación de desigualdades válidas— sin requerir el aparato poliédrico ni los
mecanismos de generación de columnas propios de los algoritmos contemporáneos. La
resolución del subproblema entero en cada iteración se delega al solver
COIN-OR CBC \parencite{pulp2011}, que internamente implementa \textit{Branch and Bound}
sobre la formulación enriquecida iterativamente con los cortes generados.

\subsection{Aplicaciones del CVRP en investigación operacional}

Desde el punto de vista aplicado, el CVRP ha servido como modelo base para problemas de
abastecimiento urbano, distribución de última milla, recolección de residuos,
planificación de rutas de combustibles y despacho a puntos de servicio
\parencite{tan2021vehicle}. En muchos de estos dominios, las versiones reales incorporan
restricciones adicionales como ventanas de tiempo, flotas heterogéneas o múltiples
depósitos; sin embargo, el CVRP sigue cumpliendo un rol metodológico central porque
permite aislar el efecto de la capacidad vehicular y de la asignación de clientes a
rutas. Por ello, continúa siendo un banco de pruebas privilegiado para comparar
algoritmos exactos sobre instancias estandarizadas y soluciones de referencia.